% Page 33 du document source
% Transcription provisoire — à relire et corriger si nécessaire.

\begin{center}
    {\Large\itshape Alban au XV\textsuperscript{ème} siècle (suite)}
\end{center}

\vspace{1.5em}

Une autre notabilité, Paulin Limozin, achète à Armand de Rames,
commandeur des hospitaliers de Rayssac, des terres, des bois et des prés,
sis à Montredon, en 1447. Et c'est ainsi qu'un jour sera fondée sur ces
terres la chapellenie de Sainte Anne. Avec les revenus de la métairie de
Négremontet, entre le rec de l'Espazet et le chemin de la Tapie à Alban
jusqu'au rec de Fangouse et les prés de la Maye, devront être assurées à
perpétuité deux messes chaque semaine au bénéfice des donateurs. Trois
siècles après, les biens de la dite chapellenie seront cédés pour 1125
francs (de l'époque), le 17 février 1792 à Marcel Cros de Carménel qui les
avait lui-même achetés à Bonnet (de Carménel). Ce dernier les avait
acquis, lors de la vente des biens nationaux (biens confisqués par les
révolutionnaires) par les administrateurs au district d'Albi.

\medskip

Il se pourrait donc que les tènements plus ou moins en copropriété des
habitants de Carménel (en patois « patus ») soient peut-être, les biens
qui assuraient cette fondation de la chapellenie Sainte Anne. La Croix
qui a été dressée en ces parages, ces derniers temps, au sommet de la côte
de Fangouse et à la limite de terrains plus ou moins en friche, rappelle
ce passé. C'est la Croix de Sainte Anne, vocable que portait la susdite
chapellenie (chapellenie Sainte Anne). Si seulement, tous les ans, chaque
habitant de Carménel se donnait la peine d'y planter une sapinette, dans
vingt ans, il y aurait aux portes d'Alban un magnifique bois… à portée de
la main… et rentable.
